\subsubsection{Quantidade de movimento angular do sistema}
\label{sec:dinAcoplada_momento_angular}

Para fins de análise física e verificação da conservação da quantidade de movimento angular,
introduz-se o cálculo explícito da quantidade de movimento angular total da espaçonave em relação à origem do referencial $\{B\}$.
A formulação de velocidades espaciais adotada na Seção~\ref{sec:dinAcoplada_energia_cinetica}, a velocidade espacial do corpo $i$ (espaçonave ou \gls{mr}) em $\{B\}$ é escrita como

\begin{equation}
\vec v_{i,\text{sp}}
=
\begin{bmatrix}
{}^{B}\!\vec v_{G_i} \\[2pt]
{}^{B}\!\vec\omega_i
\end{bmatrix}
=
J_{\text{sp},i}(\vec q\,)\,\dot{\vec q\,},
\end{equation}

onde ${}^{B}\!\vec v_{G_i}$ e ${}^{B}\!\vec\omega_i$ são, respectivamente,
a velocidade linear do centro de massa e a velocidade angular do corpo $i$
expressas em $\{B\}$, e $J_{\text{sp},i}(\vec q\,)$ é o jacobiano espacial
avaliado no centro de massa do corpo $i$, construído a partir dos graus de
liberdade de atitude da base e das juntas do manipulador
\cite{featherstone2008}.

Para cada corpo $i$, o operador de inércia espacial
$M_{i,\text{spatial}}$ associa a velocidade espacial ao momento espacial
\cite{featherstone2008}:

\begin{equation}
\vec h_{i,\text{sp}}
=
M_{i,\text{spatial}}\,\vec v_{i,\text{sp}},
\end{equation}

onde $\vec h_{i,\text{sp}}$ agrupa em um único vetor o momento linear e o momento angular do corpo $i$ em relação à origem de $\{B\}$.
A parte angular de $\vec h_{i,\text{sp}}$ é obtida selecionando-se as
três últimas componentes, por meio da matriz

\begin{equation}
S_\omega
=
\begin{bmatrix}
0 & 0 & 0 & 1 & 0 & 0 \\
0 & 0 & 0 & 0 & 1 & 0 \\
0 & 0 & 0 & 0 & 0 & 1
\end{bmatrix},
\qquad
\vec h_i
=
S_\omega\,\vec h_{i,\text{sp}},
\end{equation}

em que $\vec h_i$ é o vetor de momento angular do corpo $i$ em $\{B\}$,
em analogia à decomposição do momento espacial em componentes linear e
angular discutida em \cite{featherstone2008,lynch2017modern}.

Neste trabalho, a base é tratada como um corpo rígido adicional com operador de inércia espacial
$M_{\text{base},\text{spatial}}$ expresso em $\{B\}$,
cujos parâmetros de massa, centro de massa e tensor de inércia já incluem
a contribuição estrutural da espaçonave.
Os elos do manipulador são indexados por $i=1,\dots,n$, de modo que o
momento angular total do sistema espaçonave--manipulador em $\{B\}$ é escrito como

\begin{equation}
\vec h_{\text{tot}}
=
\sum_{i=0}^{n} \vec h_i
=
\sum_{i=0}^{n}
S_\omega\,M_{i,\text{spatial}}\,
J_{\text{sp},i}(\vec q\,)\,\dot{\vec q\,}.
\end{equation}

Essa expressão define o momento espacial total como a soma dos momentos de cada corpo rígido
em sistemas de base flutuante \cite{featherstone2008} e com formulações de sistemas espaçonave--manipulador
que utilizam a conservação de momento para descrever a dinâmica em
microgravidade.

Na ausência de torques externos, a quantidade de movimento angular total em
torno do centro de massa da espaçonave robótica é conservado, de modo que
$\vec h_{\text{tot}}$ deve permanecer constante no tempo \cite{wie2008space,featherstone2008}.